% !TEX root = ../../../main.tex
% Block card: detection.freq_stft -- "STFT spectrum" (Part I, the blocks chapter).
\begin{plainterms}[A spectrum]
A spectrum says how much of a signal's wiggle happens at each speed: slow swells at the low
frequencies, fast flicker at the high ones, like the bass and treble bars on a music
player. The block measures it on short overlapping windows and averages them. It groups
frequencies into bands that widen as the frequency rises, so the slow end, where most of
the interest lies, gets more and narrower bands.
\end{plainterms}

\cardfield{What it does}
Measures how the span's power is spread across frequencies, averaged over the whole span,
in bands spaced evenly on a logarithmic scale.

\cardfield{When to use it}
To compare the frequency content of spans or channels, or to choose a band before a
band-pass filter. It gives one spectrum for the span; it does not show how the spectrum
changes over time.

\cardfield{What goes in, what comes out}
In goes a span of one channel. Out comes one list of numbers, the average power in each
band, with each band's centre frequency; the block page draws it as a curve.

\cardfield{Settings that matter}
\emph{window\_size} is the window length in samples (default 300): longer windows resolve
slower frequencies, down to one cycle per window. \emph{hop\_size} is the step between
windows (default 150, so windows overlap by half). \emph{n\_log\_bins} is the number of
bands (default 64).

\cardfield{Watch out for}
At the slow end the bands are narrower than the window can resolve, so many of them hold no
frequency at all and read as exactly zero power: at the defaults, roughly the lowest third.
A zero there means ``no measurement'', not ``no power''. Power is in the recording's unit
squared, and a single missing sample (NaN) turns the whole spectrum into NaN. Settings
count samples, so the same numbers mean different frequencies at 1\,Hz and at 10\,Hz.
